\documentclass{situationist-preprint}

\title{Resemblance and Expression}
\msid{AES-07-b}
\msfamily{AES-07}
\msdiscipline{Aesthetics}
\mstype{Companion manuscript}
\mskeywords{musical expression; contour theory; resemblance; Goodman; asymmetry; appearance emotionalism}

\begin{document}
\maketitle

\begin{abstract}
AES-07-a argues that music expresses emotions by resembling their characteristic expressive behaviour. Resemblance has a bad reputation as a basis for anything semantic: Goodman argued that it is reflexive and symmetric, while representation and expression are neither. We argue that the objection, which is decisive against resemblance theories of pictorial representation as Goodman construed them, does not touch the contour theory. Expression in music is not resemblance simpliciter but heard resemblance, in a respect fixed by the perceptual system that recognizes emotion in behaviour, and this relation is neither reflexive nor symmetric. We state the relation precisely, show that it has the formal properties expression requires, and note what the empirical evidence on shared acoustic cues adds.
\end{abstract}

\section{Goodman's objection}

\citet{Goodman1968} argued that resemblance cannot be the basis of representation because it has the wrong formal properties. Resemblance is reflexive, since everything resembles itself to the highest degree, and symmetric, since if $a$ resembles $b$ then $b$ resembles $a$. Representation is neither: a painting does not represent itself, and a duke represented by a portrait does not represent the portrait. The same argument applies to expression. A sad passage expresses sadness; the drooping gait it resembles does not express the passage, and the passage does not express itself.

AES-07-a replies to a different objection of Goodman's, that everything resembles everything in some respect \citep{Goodman1970}, by restricting the respects to those listed in Definition~2.2: tempo, melodic contour and register, loudness, articulation and rhythm, compared with rate of movement, vocal pitch, intensity and the smoothness of movement. The present objection survives that restriction, since restricted resemblance is still reflexive and symmetric.

\section{Heard resemblance}

\begin{definition}\label{d:heard}
A passage $M$ is \emph{heard as resembling} the expressive behaviour of an emotion $E$ by a listener $L$ if $L$'s capacity to recognize $E$ in human behaviour is engaged by the dynamic profile of $M$, so that $L$ hears $M$ as having the character of $E$-behaviour without taking $M$ to be such behaviour.
\end{definition}

The definition makes explicit what Principle~P2 of AES-07-a requires, that expressiveness be heard in the music. It corresponds to what \citet{Kivy1980} and \citet{SDavies1994} described as the perception of emotion characteristics in appearances.

\begin{proposition}\label{p:asym}
Heard resemblance in the sense of Definition~\ref{d:heard} is neither reflexive nor symmetric.
\end{proposition}

\begin{proof}
It is not reflexive: the relation holds between a passage and a kind of behaviour, and a passage is not a kind of behaviour, so the relation has no instance of the form $R(x,x)$. It is not symmetric: the capacity engaged is a capacity to recognize emotion in behaviour. When a listener hears a passage as resembling a sad gait, her recognition of sadness is engaged by the passage; when she sees a sad gait, her recognition of sadness is engaged by the gait, and she does not thereby hear any passage in it. The relation runs from the stimulus to the recognitional capacity whose proper domain is behaviour, and only one relatum lies outside that domain.
\end{proof}

\begin{proposition}\label{p:expr}
If a passage contour-resembles the expressive behaviour of $E$ in the sense of Definition~2.2 of AES-07-a, and the resemblance is heard in the sense of Definition~\ref{d:heard} by competent listeners, the passage is expressive of $E$ in the sense of Theorem~3.4 of AES-07-a, and the relation has the formal properties Goodman's objection demands.
\end{proposition}

\begin{proof}
Contour resemblance provides the respect; hearing it engages the recognitional capacity that makes the resemblance expressive rather than merely structural; and by Proposition~\ref{p:asym} the composite relation is irreflexive and asymmetric. The same holds of a weeping willow seen as drooping, which AES-07-a, in its reply to its second objection, grants is expressive of sadness in this sense; what distinguishes music is that it is made and attended to for this appearance, which does not alter the formal properties of the relation.
\end{proof}

The empirical evidence fits the proposal. \citet{JuslinLaukka2003} found that the acoustic cues by which listeners identify emotions in speech, such as rate, intensity, pitch level and variability, and the sharpness of onsets, are largely the same as those by which they identify emotions in music performance. This is the overlap that a single recognitional capacity engaged by both would predict.

\begin{objection}[Projection]
If expression depends on the listener's recognitional capacity, expression is in the listener and not in the music.
\end{objection}

\begin{reply}
The capacity is the listener's; what it responds to is in the music. A passage is expressive of sadness if it is apt to engage competent listeners' recognition of sadness through its dynamic profile, and whether it is so apt is fixed by that profile. This is a response-dependent property of the kind AES-02 identifies for beauty, and response-dependence does not locate a property in the observer.
\end{reply}

\section{Conclusion}

Goodman's argument shows that resemblance alone cannot be expression. It does not show that heard resemblance, in respects fixed by the recognition of emotion in behaviour, cannot be. That relation is irreflexive and asymmetric, it explains why a gait does not express the passage that resembles it, and it is the relation the contour theory of AES-07-a requires.

\begin{thebibliography}{99}\small
\bibitem[Davies(1994)]{SDavies1994} Davies, S. (1994). \emph{Musical Meaning and Expression}. Cornell University Press.
\bibitem[Goodman(1968)]{Goodman1968} Goodman, N. (1968). \emph{Languages of Art: An Approach to a Theory of Symbols}. Bobbs-Merrill.
\bibitem[Goodman(1970)]{Goodman1970} Goodman, N. (1970). Seven strictures on similarity. In L. Foster and J.~W. Swanson (Eds.), \emph{Experience and Theory}. University of Massachusetts Press.
\bibitem[Juslin and Laukka(2003)]{JuslinLaukka2003} Juslin, P.~N. and Laukka, P. (2003). Communication of emotions in vocal expression and music performance: Different channels, same code? \emph{Psychological Bulletin}, 129(5), 770--814.
\bibitem[Kivy(1980)]{Kivy1980} Kivy, P. (1980). \emph{The Corded Shell: Reflections on Musical Expression}. Princeton University Press.
\end{thebibliography}

\end{document}
